% ch11_b §11.3–11.5 (书页 386–393 / p400–p407)
%--A-TAIL--
%JOIN%（在 BCS 理论中，这个量记作 $N(0)$，即费米能级处\emph{每个自旋}的态密度）。把式 (11.6) 代入式 (11.20)，并取 $w\approx k\Theta$，就应当得到
\begin{align*}
\Delta_{0} &\approx 2k\Theta\exp\{-2/\mathcal{N}(\mathcal{E}_{F})\,|V|\}\\
&\approx 2k\Theta\exp\{-8\mathcal{E}_{F}^{2}/9\,|\mathcal{U}|^{2}\}.
\tag{11.22}
\end{align*}
如果我们假定 $\Delta_{0}$ 是超导转变能量的某种量度，并令它等于 $kT_{c}$，就会发现转变温度 $T_{c}$ 应当介于 Debye 温度的 1/10 与 1/10,000 之间，视电子–离子相互作用势 $|\mathcal{U}|$ 的大小而定。实际观察到的正是如此——我们现在可以理解，为什么 $T_{c}$ 与 $\Theta$ 相差得如此悬殊。

由这个论证还可得出的另一个效应是\emph{同位素效应}（isotope effect）。如果考虑改变每个离子的质量 $M$ 而不改变原子间的作用力，那么我们预期声速按 $M^{-1/2}$ 变化。于是我们预期 $\Theta$ 也按同样的方式变化，从而
\begin{equation*}
T_{c}\propto M^{-1/2},
\tag{11.23}
\end{equation*}
这在许多情形中与观察结果大体相符。

关于 Cooper 对还有一点：如果两个电子除动量相反之外自旋也相反，它们将束缚得更牢。这一点从我们刚才给出的论证看并不明显，而是在计算吸引势 (11.2) 时显露出来的。§11.1 的论证假定我们处理的是可分辨粒子，即自旋相反的电子。倘若两个电子自旋相同，我们就必须援引 Pauli 原理，在计算有效相互作用时使用反对称化的波函数。“交换”项的计入会削弱吸引力，因为这时我们必须设法阻止两个电子进入同一状态或占据同一位置。Cooper 对的单态就是这个准分子的最低态。

\section{超导基态}

上述结果提示了超导态的两个重要特征：它是一个无法用微扰展开中的有限多项来达到的态，因为函数 (11.20) 在 $V=0$ 附近不是 $V$ 的解析函数；它是一个电子以某种方式配成对的态，$\mathbf{k}$ 与 $-\mathbf{k}$ 相配。这个态的构造最初是由 Bardeen、Cooper 和 Schrieffer 所完成的。为了演示他们的解，我们将采用 Bogoliubov 的方法。

完全可以避开有关波函数的一切讨论，而把整个问题用算符语言表述出来，这是非常方便的。我们需要表述这样一个事实：我们有 $N$ 个电子，其能量主要是动能，处在由矢量 $\mathbf{k}$ 和自旋指标 $\sigma$ 标记的态上，并且它们通过势 $V(\mathbf{K})$ 发生相互作用。我们定义湮没算符和产生算符 $b_{\mathbf{k}\sigma}$、$b^{*}_{\mathbf{k}\sigma}$，它们类似于式 (2.129) 和 (10.109) 中定义的算符 $a_{\mathbf{q}}$、$a^{*}_{\mathbf{q}}$，只不过我们现在处理的是费米子。这意味着它们必须服从\emph{反对易关系}（anticommutation relations）
\begin{align*}
[b_{\mathbf{k}\sigma},\,b^{*}_{\mathbf{k}'\sigma'}]_{+} &\equiv b_{\mathbf{k}\sigma}b^{*}_{\mathbf{k}'\sigma'}+b^{*}_{\mathbf{k}'\sigma'}b_{\mathbf{k}\sigma}\\
&=\delta_{\mathbf{k}\mathbf{k}'}\,\delta_{\sigma\sigma'}.
\tag{11.24}
\end{align*}
这个定义的结果是，算符
\begin{equation*}
n_{\mathbf{k}\sigma}=b^{*}_{\mathbf{k}\sigma}b_{\mathbf{k}\sigma}
\tag{11.25}
\end{equation*}
量度态 $(\mathbf{k},\sigma)$ 的占据数，并且如 Pauli 原理所要求的那样，本征值取 0 和 1。当式 (11.24) 中的 $\mathbf{k}$ 与 $\mathbf{k}'$、或 $\sigma$ 与 $\sigma'$ 不相同时，反对易子体现的是反对称原理：交换两个态的标记将使波函数反号。

我们这个系统的 Hamilton 量可以写成
\begin{equation*}
\mathcal{H}=\sum_{\mathbf{k}\sigma}\xi(\mathbf{k})\,b^{*}_{\mathbf{k}\sigma}b_{\mathbf{k}\sigma}+\sum_{\mathbf{k}_1,\,\mathbf{k}_2,\,\mathbf{K}}V(\mathbf{K})\,b^{*}_{\mathbf{k}_2-\mathbf{K},\downarrow}\,b^{*}_{\mathbf{k}_1+\mathbf{K},\uparrow}\,b_{\mathbf{k}_1,\uparrow}\,b_{\mathbf{k}_2,\downarrow}.
\tag{11.26}
\end{equation*}
第一项代表电子的动能，像式 (11.12) 中那样，从费米能级算起。为了方便，常常假定这个能量可以变动——即电子总数 $n$ 并非严格固定——办法是引进电子气的化学势 $\zeta$，并把第一项写成
\begin{align*}
\mathcal{H}_{0}&=\sum_{\mathbf{k}\sigma}n_{\mathbf{k}\sigma}\,\xi(\mathbf{k})\\
&=\sum_{\mathbf{k}\sigma}n_{\mathbf{k}\sigma}\{\mathcal{E}(\mathbf{k})-\zeta\}\\
&=\sum_{\mathbf{k}\sigma}n_{\mathbf{k}\sigma}\,\mathcal{E}(\mathbf{k})-n\zeta.
\tag{11.27}
\end{align*}
式 (11.26) 中的第二项代表由相互作用 (11.2) 所定义的散射，其效果是把 $\mathbf{k}_1$ 和 $\mathbf{k}_2$ 中的电子“湮没”掉，然后再把它们“重新产生”在 $\mathbf{k}_1+\mathbf{K}$ 和 $\mathbf{k}_2-\mathbf{K}$ 处，矩阵元为 $V(\mathbf{K})$。这里我们假定只有自旋相反的电子之间才有相互作用。

Bogoliubov 方法是对算符集 $b_{\mathbf{k}\sigma}$ 与 $b^{*}_{\mathbf{k}\sigma}$ 作一个正则变换，把它们变成具有同样对易关系的一套新的湮没算符和产生算符。特别地，这些新算符必须显示出某种配对性质——正如 Cooper 效应所提示的，它们必须把态 $(\mathbf{k},\uparrow)$ 与态 $(-\mathbf{k},\downarrow)$ 联结起来。

为简单起见，让我们略去自旋指标，写成
\begin{equation*}
\left.\begin{array}{ll}
\beta^{*}_{\mathbf{k}}=u_{\mathbf{k}}b^{*}_{\mathbf{k}}-v_{\mathbf{k}}b_{-\mathbf{k}}, &\quad \beta_{\mathbf{k}}=u_{\mathbf{k}}b_{\mathbf{k}}-v_{\mathbf{k}}b^{*}_{-\mathbf{k}},\\
\beta^{*}_{-\mathbf{k}}=u_{\mathbf{k}}b^{*}_{-\mathbf{k}}+v_{\mathbf{k}}b_{\mathbf{k}}, &\quad \beta_{-\mathbf{k}}=u_{\mathbf{k}}b_{-\mathbf{k}}+v_{\mathbf{k}}b^{*}_{\mathbf{k}}.
\end{array}\right\}
\tag{11.28}
\end{equation*}
（译注：原书 (11.28) 第一式排印为“$\beta_{\mathbf{k}}\ \ u_{\mathbf{k}}b^{*}_{\mathbf{k}}-v_{\mathbf{k}}b_{-\mathbf{k}}$”，漏印等号与 $\beta$ 上的星号；按变换关系此处应为第二式的 Hermite 共轭，译文已补正。）

如果
\begin{equation*}
u^{2}_{\mathbf{k}}+v^{2}_{\mathbf{k}}=1
\tag{11.29}
\end{equation*}
则算符 $\beta_{\mathbf{k}}$、$\beta^{*}_{\mathbf{k}}$ 就会具有费米子应有的反对易关系。
这个变换确实与反铁磁自旋波的变换 (10.130) 本质上相同，只不过现在我们是在 $\mathbf{k}$ 与 $-\mathbf{k}$ 的空间中作一个“真实转动”，于是
\begin{equation*}
u_{\mathbf{k}}=\cos\theta_{\mathbf{k}},\quad v_{\mathbf{k}}=\sin\theta_{\mathbf{k}}
\tag{11.30}
\end{equation*}
满足式 (11.29)。这是因为我们现在处理的是费米子，而不是玻色子。

当我们从式 (11.28) 解出原来的算符并代入 Hamilton 量 (11.26) 时，会得到形形色色的项，其中含有新算符 $\beta_{\mathbf{k}}$、$\beta^{*}_{\mathbf{k}}$ 的乘积。现在，只要某一项中含有湮没算符 $b_{\mathbf{k}}$，我们就可以利用反对易子把它移到右边，例如
\begin{equation*}
\beta_{\mathbf{k}}\beta^{*}_{\mathbf{k}'}=\delta_{\mathbf{k}\mathbf{k}'}-\beta^{*}_{\mathbf{k}'}\beta_{\mathbf{k}}.
\tag{11.31}
\end{equation*}
假定我们的超导态就是这些新算符的“真空”$|0\rangle$。按定义，真空中不存在任何能被湮没的对象，因此
\begin{equation*}
\beta_{\mathbf{k}}|0\rangle=0,
\tag{11.32}
\end{equation*}
对所有 $\mathbf{k}$ 成立。Hamilton 量中的这类项对基态 $|0\rangle$ 的能量没有贡献；于是基态是 $\mathcal{H}$ 中这一部分的本征态。

但是还剩下一些项不能用这种方式消去——即只含有像 $\beta^{*}_{\mathbf{k}}$ 这样的产生算符的项。它们来自 (11.26) 的动能部分，也来自相互作用项的约化——例如当式 (11.31) 中的 $\mathbf{k}$ 与 $\mathbf{k}'$ 恰好相等时。我们发现，Hamilton 量可以写成
\begin{align*}
\mathcal{H}&=2\sum_{\mathbf{k}}\xi(\mathbf{k})\,v^{2}_{\mathbf{k}}+\sum_{\mathbf{k},\,\mathbf{K}}V(\mathbf{K})\,u_{\mathbf{k}}v_{\mathbf{k}}u_{\mathbf{k}+\mathbf{K}}v_{\mathbf{k}+\mathbf{K}}\\
&\quad+\sum_{\mathbf{k}}\bigl\{2\xi(\mathbf{k})\,u_{\mathbf{k}}v_{\mathbf{k}}+(u^{2}_{\mathbf{k}}-v^{2}_{\mathbf{k}})\sum_{\mathbf{K}}V(\mathbf{K})\,u_{\mathbf{k}+\mathbf{K}}v_{\mathbf{k}+\mathbf{K}}\bigr\}\beta^{*}_{\mathbf{k}}\beta^{*}_{-\mathbf{k}}.
\tag{11.33}
\end{align*}
前面的这些项将给出我们基态的能量——只要我们能把最后那个求和中的产生算符消去。

这跟我们在 (10.132) 中遇到的情况非常相似——我们采用同样的手法。我们这样选取变换系数 $u_{\mathbf{k}}$、$v_{\mathbf{k}}$，使每个“危险”项都等于零。令
\begin{equation*}
2\xi(\mathbf{k})\,u_{\mathbf{k}}v_{\mathbf{k}}+(u^{2}_{\mathbf{k}}-v^{2}_{\mathbf{k}})\sum_{\mathbf{K}}V(\mathbf{K})\,u_{\mathbf{k}+\mathbf{K}}v_{\mathbf{k}+\mathbf{K}}=0.
\tag{11.34}
\end{equation*}
一般说来，这是关于未知函数 $u_{\mathbf{k}}$ 和 $v_{\mathbf{k}}$ 的一个复杂的积分方程，而且它们还须满足条件 (11.29)。但让我们采用 (11.8) 和 (11.9) 的假定，把吸引势 $V(\mathbf{K})$ 简化为在能量范围 $\pm w$ 之内取常数值 $V$、在此范围之外为零。我们把式 (11.34) 中对 $\mathbf{K}$ 的求和记作一个参量
\begin{equation*}
\Delta_{0}=-V\sum_{-w}^{w}u_{\mathbf{k}+\mathbf{K}}v_{\mathbf{k}+\mathbf{K}},
\tag{11.35}
\end{equation*}
然后联立求解 (11.34) 与 (11.29)。结果是
\begin{equation*}
u^{2}_{\mathbf{k}}=\frac{1}{2}\biggl[1+\frac{\xi(\mathbf{k})}{\sqrt{\{\Delta^{2}_{0}+\xi^{2}(\mathbf{k})\}}}\biggr],\quad v^{2}_{\mathbf{k}}=\frac{1}{2}\biggl[1-\frac{\xi(\mathbf{k})}{\sqrt{\{\Delta^{2}_{0}+\xi^{2}(\mathbf{k})\}}}\biggr],
\tag{11.36}
\end{equation*}
把它代回 (11.35)，便得到
\begin{align*}
1&=-\frac{1}{2}V\sum_{-w}^{w}\{\Delta^{2}_{0}+\xi^{2}(\mathbf{k})\}^{-\frac{1}{2}}\\
&=-\frac{1}{2}\,\mathcal{N}(\mathcal{E}_{F})\,V\ln\frac{2w}{\Delta_{0}},
\tag{11.37}
\end{align*}
这是在费米面附近对 $\xi$ 的范围积分而来的（记住自旋只计一种）。这与 (11.19) 完全相同。参量 $\Delta_{0}$ 恰恰就是一个 Cooper 对的结合能，与 (11.20) 或 (11.22) 中的一样。

我们可以把这个结果代回 (11.33) 的其余部分——现在其中所有的算符都已消去。我们得到
\begin{equation*}
\mathcal{E}_{0}=2\sum_{\xi<0}\xi(\mathbf{k})-\frac{1}{2}\Delta^{2}_{0}\sum_{\mathbf{k}}\{\Delta^{2}_{0}+\xi^{2}(\mathbf{k})\}^{-\frac{1}{2}}.
\tag{11.38}
\end{equation*}
于是，态 $|0\rangle$ 是 $\mathcal{H}$ 的一个本征态——但其能量低于电子填满的费米球，后者的能量本来正好是各被占据态的 $\xi(\mathbf{k})$ 值之和，如 (11.27) 那样。还可以证明，对于电子总数的期望值 $n$ 的变动而言，这个态使能量取极小，正如 (11.27) 所要求的。

\section{准粒子与能隙}

超导态 $|0\rangle$ 的本性是什么？按定义，它是算符 $\beta^{*}_{\mathbf{k}}$、$\beta_{\mathbf{k}}$ 的“真空”。这些算符能够产生和消灭真空的激发，正如原来的算符 $b^{*}_{\mathbf{k}}$ 和 $b_{\mathbf{k}}$ 能够“产生”和“消灭”电子一样。这样一种激发看上去会是什么样子？按照 (11.28)，$\beta^{*}_{\mathbf{k}}$ 的效果是以振幅 $u_{\mathbf{k}}$“产生”一个处在态 $\mathbf{k}$ 中的电子，同时又以振幅 $v_{\mathbf{k}}$“消灭”一个处在态 $-\mathbf{k}$ 中的电子。

这些系数由 (11.36) 给出。假定 $\mathbf{k}$ 远在费米面之上，使得 $\xi(\mathbf{k})\gg\Delta_{0}$，那么 $u^{2}_{\mathbf{k}}\approx 1$，而 $v^{2}_{\mathbf{k}}\approx 0$。我们的激发看上去就会像一个普通的电子。
%FIG{203}: {准粒子。}
但是现在，当我们逼近费米能级时，$\xi(\mathbf{k})$ 趋于零，$b^{*}_{\mathbf{k}}$ 的振幅将减小，而 $b_{-\mathbf{k}}$ 的振幅将增大。当 $\xi(\mathbf{k})$ 为负时，$v^{2}_{\mathbf{k}}$ 的值将大于 $\frac{1}{2}$；当我们远在费米能级之下时，$v^{2}_{\mathbf{k}}\approx 1$ 而 $u^{2}_{\mathbf{k}}\approx 0$。这时我们的激发相应于在 $-\mathbf{k}$ 处一个电子的“毁灭”，也就是说，它是本已被占据的诸态中的一个“空穴”。

因此，超导态的激发是颇为奇特的\emph{准粒子}（quasi-particle），当它们通过费米能级时，会从“电子”变成“空穴”。在 $\Delta_{0}$ 附近的能量范围内，每个准粒子都是 $\mathbf{k}$ 处的一个电子与 $-\mathbf{k}$ 处的一个空穴的混合体。Cooper 对本身不再出现，但是动量相反（当然还有自旋相反）的电子之间的关联是显而易见的。实际上，我们是在坚持：配对态中各波函数的相位之间应当存在一种关系。

激发的能量也可以计算。我们在 (11.26) 的展开式中寻找含下列算符
\begin{equation*}
n'_{\mathbf{k}}=\beta^{*}_{\mathbf{k}}\beta_{\mathbf{k}}
\tag{11.39}
\end{equation*}
的各项的系数，因为它量度的是在此模式中被激发的准粒子的数目。结果是
\begin{align*}
\epsilon(\mathbf{k})&=\xi(\mathbf{k})\,(u^{2}_{\mathbf{k}}-v^{2}_{\mathbf{k}})-2u_{\mathbf{k}}v_{\mathbf{k}}\sum_{\mathbf{K}}V(\mathbf{K})\,u_{\mathbf{k}+\mathbf{K}}v_{\mathbf{k}+\mathbf{K}}\\
&=\sqrt{\{\Delta^{2}_{0}+\xi^{2}(\mathbf{k})\}}
\tag{11.40}
\end{align*}
（对我们的简化相互作用，用了 (11.35) 和 (11.36)）。
%FIG{204}: {准粒子的能量。}
这一点非常重要，因为它证明了在超导基态之上存在\emph{能隙}（energy gap）$\Delta_{0}$。当 $\xi(\mathbf{k})\gg\Delta_{0}$ 时，$\epsilon(\mathbf{k})\approx\xi(\mathbf{k})$——准粒子的能量就是把一个普通电子激发到费米能级之上所需的能量。但当 $\mathbf{k}$ 趋近费米面（那里 $\xi(\mathbf{k})$ 为零）时，准粒子的能量趋于常数值 $\Delta_{0}$。然后，随着 $\mathbf{k}$ 再进一步减小，$\epsilon(\mathbf{k})$ 又重新增大，直至几乎等于 $|\xi(\mathbf{k})|$。这当然就是把一个电子从负能量态 $\xi(\mathbf{k})$ 中移出所需的能量，正如我们根据对准粒子算符的解释所预期的那样。
因此，超导态是一种“凝聚”态，其含义是：要使整个系统产生一个激发态，需要有限的能量 $\Delta_{0}$。在某些方面，它有点像半导体——在半导体中，要想把一个电子送入导带，也必须越过一有限的带隙。这个类比有助于粗略地解释超导体的某些特殊性质，例如电磁吸收（参见 §8.5）和隧穿（§6.8）。

\section{能隙随温度的变化}

基态只适用于 $T=0$ 的情形。在有限温度下，我们预期会有准粒子按通常的 费米–狄拉克 函数被激发：
\begin{equation*}
f_{\mathbf{k}}=\frac{1}{e^{\epsilon(\mathbf{k})/kT}+1}.
\tag{11.41}
\end{equation*}
在正常态中，这类激发无论在费米能级之上还是之下，都彼此独立。在超导态中，它们倾向于相互合作地发生作用，并破坏能隙。

让我们假定所处理的是态 $|f_{\mathbf{k}}\rangle$，其中第 $\mathbf{k}$ 个态中准粒子的平均数目由 (11.41) 给出。湮没算符 $b_{\mathbf{k}}$ 作用在这个态上的效果不再是零。平均而言，我们可以把 (11.25) 换成
\begin{equation*}
b^{*}_{\mathbf{k}}b_{\mathbf{k}}|f_{\mathbf{k}}\rangle=f_{\mathbf{k}}|f_{\mathbf{k}}\rangle.
\tag{11.42}
\end{equation*}
把 Hamilton 量约化成对角形式的做法，当应用于态 $|0\rangle$ 时，依靠的是 (11.32)。如果现在对态 $|f_{\mathbf{k}}\rangle$ 试着做同样的事情，就会多出一些额外的项；它们来自像 (11.31) 那样把算符乘积约化成标准次序的过程，并乘在 (11.33) 中那些“危险”乘积 $\beta^{*}_{\mathbf{k}}\beta^{*}_{-\mathbf{k}}$ 上。为了从 Hamilton 量中去掉这些“产生电子对”的部分，我们的条件 (11.34) 必须改成为
\begin{equation*}
2\xi(\mathbf{k})\,u_{\mathbf{k}}v_{\mathbf{k}}+(u^{2}_{\mathbf{k}}-v^{2}_{\mathbf{k}})\sum_{\mathbf{K}}V(\mathbf{K})\,u_{\mathbf{k}+\mathbf{K}}v_{\mathbf{k}+\mathbf{K}}(1-2f_{\mathbf{k}+\mathbf{K}})=0
\tag{11.43}
\end{equation*}
（这里计及与每个 $\mathbf{k}$ 相联系的两个可能的自旋态）。

对于简化的相互作用 $V$，这个积分方程可以这样求解：把 (11.35) 和 (11.36) 中的 $\Delta_{0}$ 换成
\begin{equation*}
\Delta=-V\sum_{-w}^{w}u_{\mathbf{k}+\mathbf{K}}v_{\mathbf{k}+\mathbf{K}}(1-2f_{\mathbf{k}+\mathbf{K}}).
\tag{11.44}
\end{equation*}
代替 (11.37)，利用 (11.41)，我们便得到
\begin{align*}
1&=-\frac{1}{2}V\sum_{-w}^{w}\frac{1}{\epsilon(\mathbf{k})}\tanh\frac{\epsilon(\mathbf{k})}{2kT}\\
&=-\frac{1}{2}\,\mathcal{N}(\mathcal{E}_{F})\,V\int_{-w}^{w}\frac{1}{\epsilon}\tanh\frac{\epsilon}{2kT}\,d\xi,
\tag{11.45}
\end{align*}
其中 $\epsilon$ 和以前一样，是一个激发的能量，
\begin{equation*}
\epsilon(\mathbf{k})=\sqrt{\{\Delta^{2}+\xi^{2}(\mathbf{k})\}}.
\tag{11.46}
\end{equation*}
式 (11.45) 与式 (11.46) 这两个方程给出 $\Delta$ 与 $T$ 之间的一个隐含关系。这个关系相当杂乱，但要点在于：当 $T$ 增大时，$\Delta$ 从 $T=0$ 处的 $\Delta_{0}$ 逐渐减小。换言之，能隙 $\Delta$ 随温度升高而减小，并在一个完全确定的温度 $T_{c}$ 处闭合为零；由 (11.46)，在该处 $\epsilon(\mathbf{k})=|\xi(\mathbf{k})|$。于是，我们可以通过积分
\begin{equation*}
\frac{2}{\mathcal{N}(\mathcal{E}_{F})\,V}=2\int_{0}^{w}\frac{1}{\epsilon}\tanh\left(\frac{\epsilon}{2kT_{c}}\right)d\epsilon,
\tag{11.47}
\end{equation*}
求出 $T_{c}$，其解为
\begin{equation*}
\Delta_{0}\approx 1.76\,kT_{c},
\tag{11.48}
\end{equation*}
%FIG{205}: {能隙随温度的变化。}
式中的 $\Delta_{0}$ 由 (11.22) 表示，即 $T=0$ 时的能隙。这样一来，§11.2 中由 Cooper 对的结合能来估计 $T_{c}$ 的讨论就得到了印证。

在 $T_{c}$ 以上，方程没有解。恰好低于 $T_{c}$ 时，能隙从零陡峭地升起，其变化规律为
\begin{equation*}
\Delta(T)\approx 3.2\,kT_{c}\left(1-\frac{T}{T_{c}}\right)^{\frac{1}{2}},
\tag{11.49}
\end{equation*}
这可以通过在 (11.47) 之解的邻域内寻求 (11.45) 与 (11.46) 的近似解来验证。在更低的温度下，它增大得较为缓慢，在大约 $T\sim\frac{1}{2}T_{c}$ 以下逐渐变平而趋于 $\Delta_{0}$。

从对我们的准粒子系统的能量所作的仔细计算——其中计及能隙随温度
%（句子未完，下接书页 394 / p408 页首 "ture, one can write down formulae for the specific heat."——应由下一部分的 B-TAIL 以 %JOIN% 续接"的变化，便可以写下比热的公式。……"；式 (11.50)、(11.51)、图 206 及 §11.6 均在 p408 起，归下一部分。）
